\section{Related work}

\subsection{World models and transition backbones}

Contemporary world models span recurrent stochastic state-space models, discrete autoregressive transformers, implicit latent dynamics, and newer sequence architectures. DreamerV3 couples recurrent latent dynamics with normalization and imagined behavior learning, while TD-MPC2 uses decoder-free latent prediction and local trajectory optimization~\cite{hafner2025_mastering_diverse_control_tasks_through,hansen2023_td_mpc2_scalable_robust_world}. IRIS and STORM demonstrate transformer-based world models in the Atari-100k protocol~\cite{micheli2022_transformers_are_sample_efficient_world,zhang2023_storm_efficient_stochastic_transformer_based}. UniZero, HarmonyDream, and curious replay explore generalized planning, multitask harmonization, and adaptation through replay~\cite{pu2024_unizero_generalized_and_efficient_planning,ma2023_harmonydream_task_harmonization_inside_world,kauvar2023_curious_replay_for_model_based}. Scaling and regularization remain active design dimensions~\cite{nauman2024_bigger_regularized_optimistic_scaling_for}.

Applications in driving and manipulation add realistic observation and planning demands: Think2Drive and Dream-to-Drive use latent predictive models in CARLA-like settings, DINO-WM plans over pretrained visual features, and recent reviews organize learned dynamics for robotic manipulation~\cite{li2024_think2drive_efficient_reinforcement_learning_by,gao2024_dream_to_drive_with_predictive,zhou2024_dino_wm_world_models_on,ai2025_a_review_of_learning_based}. QT-TDM combines transformer dynamics with autoregressive value prediction~\cite{kotb2024_qt_tdm_planning_with_transformer}. GLAM separates global and local variation using Mamba-style sequence modules~\cite{he2025_glam_global_local_variation_awareness}. For the continuous-state forecasting task studied here, we implement compact, approximately parameter-matched adaptations of recurrent stochastic, SimNorm, and transformer transition families under a common evaluation protocol.

Recent work explores additional physical structure. The 2025 Prismatic World Model uses a mixture of experts to represent compositional hybrid modes, and the 2026 HaM-World preprint combines a soft Hamiltonian subspace, residual dynamics, and selective memory~\cite{mingwei2025_prismatic_world_model_learning_compositional,tang2026_ham_world_soft_hamiltonian_world}. Another 2026 preprint argues that structural assumptions can certify a useful prediction horizon~\cite{wang2026_scale_buys_interpolation_structure_buys}. These non-archival works provide conceptual comparisons: their mode assignment, Hamiltonian structure, and equivariant certificates address transition structure, whereas CFR-WM uses a fixed mask-exact complementary update to initialize a probabilistic residual rollout. Broader surveys place these choices within the representation, dynamics, planning, and agent-learning components of world models~\cite{ding2025_understanding_world_or_predicting_future}.

\subsection{Physics-guided dynamical learning}

Physics-guided learning can constrain outputs, embed equations, combine solvers with neural corrections, or assimilate uncertain observations. Reviews emphasize that no single integration pattern dominates across scientific problems~\cite{cuomo2022_scientific_machine_learning_through_physics,bradley2022_perspectives_on_the_integration_between,yu2024_learning_dynamical_systems_from_data}. Differentiable modeling exposes mechanistic components to gradient-based learning, while causal or horizon-aware optimization is important when temporal residuals compound~\cite{shen2023_differentiable_modelling_to_unify_machine,wang2024_respecting_causality_for_training_physics,list2022_learned_turbulence_modelling_with_differentiable}.

Several works are especially close to our transition setting. Physics-guided neural networks have been evaluated on nonlinear dynamics by injecting partial physics into a learned architecture~\cite{robinson2022_physics_guided_neural_networks_for}. Physics-guided Deep Markov Models integrate state-space physics with probabilistic latent inference~\cite{liu2022_physics_guided_deep_markov_models}. KalmanNet learns Kalman-style state updates when dynamics are only partially known~\cite{revach2022_kalmannet}. Hybrid online models have combined physical and data-driven components for electric vehicles, while related dynamic-state estimators couple both sources in integrated energy systems~\cite{chen2022_hybrid_physics_data_driven_online,chen2023_hybrid_physics_and_data_driven}. Machine learning with data assimilation further connects dynamical estimation and uncertainty quantification~\cite{cheng2023_machine_learning_with_data_assimilation}. Physics-guided lake-temperature simulation illustrates how scientific priors can be incorporated into a specific dynamical application~\cite{jia2021_physics_guided_machine_learning_for}.

CFR-WM couples mask-exact context correction with a multi-step probabilistic physics-residual world model. Its measurement gain is fixed before confirmation, making the forecast origin an explicit design choice. Gain sensitivities evaluate how that choice affects subsequent free-rollout accuracy.

\subsection{Predictive uncertainty and calibration}

Deep-learning uncertainty can represent data noise, model uncertainty, or their approximation, and a low negative log likelihood need not imply calibrated intervals~\cite{abdar2021_a_review_of_uncertainty_quantification}. Sequence-model studies in rainfall--runoff prediction provide precedents for evaluating coverage and interval width together~\cite{klotz2022_uncertainty_estimation_with_deep_learning}. Bayesian calibration work explicitly distinguishes model-form uncertainty in dynamical systems~\cite{ramancha2022_accounting_for_model_form_uncertainty}, while automated calibration has been studied for wave downscaling~\cite{alonso2021_automatic_calibration_and_uncertainty_quantification}. Conformal prediction offers a distribution-light route to uncertainty sets for dynamic biological systems~\cite{portela2025_conformal_prediction_for_uncertainty_quantification}. For the controlled benchmark, we assess the Gaussian forecasts through NLL, empirical 90\% coverage, mean interval width, and calibration error over nine central-coverage levels.
